\chapter{Statistics and Linear Algebra}

\section{Regression Models}
\noindent\textbf{Key Ideas and Methods.}\enspace
For regressions with an intercept, center the variables before working with covariances or orthogonal projections. When adding a regressor, residualize both the response and the new column against the existing design. This reduces the added direction to a one-dimensional problem. Check full rank and a nonzero residual direction before dividing by a variance or a squared norm.

Measurement error and selection change the relevant covariance structure. Specify the observed variables and the conditional population before comparing linear prediction with prediction by a conditional mean. An instrument must correlate with the latent signal and be uncorrelated with the relevant errors. Projection identities are algebraic; an exact finite-sample $t$ test additionally requires a correctly specified model with independent, homoscedastic normal errors.
\par\smallskip



\problem{A01}{Measurement Error and Instruments}{Population regression; identification}
Let $T,U,V,\varepsilon$ be mutually independent, centered random variables with finite variances, where $s=\Var(T)>0$ and $u=\Var(U)>0$. The observed variables are
\[
 W=T+U,\qquad Z=T+V,\qquad Y=\beta T+\varepsilon.
\]
Find the population regression slope of $Y$ on $W$. Can $\beta$ be identified from the covariances of $(W,Z,Y)$ alone? If $k$ readings of the same $T$, with independent identically distributed measurement errors, are averaged, how does the ordinary regression bias change? All regressions include an intercept.
\solution Independence eliminates the cross-covariances, giving
\[
 \Var(W)=s+u,\quad \Cov(W,Y)=\beta s,
 \quad \Cov(Z,W)=s,\quad \Cov(Z,Y)=\beta s.
\]
The ordinary slope is $\beta s/(s+u)$, which is attenuated toward zero. Unless $\beta=0$, it differs from the coefficient of the latent variable. Since $s>0$, the instrumental-variable ratio is well defined and identifies
\[
 \boxed{\beta=\frac{\Cov(Z,Y)}{\Cov(Z,W)}}.
\]
The instrument must correlate with the latent regressor and be uncorrelated with both the measurement error and the response noise.
For readings $W_j=T+U_j$, suppose the $U_j$ are mutually independent, independent of $(T,\varepsilon)$, and have variance $u$. Then $\bar W=T+\bar U$ has measurement-error variance $u/k$, so
\[
 \beta_k=\frac{\beta s}{s+u/k},\qquad
 \beta_k-\beta=-\frac{\beta u}{ks+u}\longrightarrow0.
\]
Averaging reduces measurement noise; it does not create $k$ independent realizations of the latent variable.
\insight Computing a projection coefficient and identifying a structural coefficient are different tasks. List the observable covariances and look for a ratio that eliminates the unknown variance.

\problem{A02}{The Gain from an Extra Regressor}{Orthogonal projection; partial determination}
Let the fixed design $X\in\R^{n\times k}$ have full column rank and contain an intercept. Add a column $z\notin\operatorname{col}(X)$, with $n>k+1$. Let $S_0$ be the residual sum of squares using $X$, and $S_1>0$ that using $(X,z)$. Write $M=I-X(X^\top X)^{-1}X^\top$ and $r=Mz$. Find the added coefficient and $S_0-S_1$, and relate the partial coefficient of determination to the usual $t$ statistic. State assumptions for an exact $t$ test.
\solution For a fixed added coefficient $b$, minimizing over the old coefficients leaves $\|My-br\|^2$. Since $r\ne0$, completing the square gives
\[
 \widehat b=\frac{r^\top y}{r^\top r},\qquad
 S_0-S_1=\frac{(r^\top y)^2}{r^\top r}=:D.
\]
The decrease is nonnegative and is zero exactly when $r^\top y=0$. It measures alignment with the old residual, rather than correlation with the raw response.
Put $\nu=n-k-1$. With $\widehat\sigma^2=S_1/\nu$, the standard error and squared statistic are
\[
 \operatorname{se}(\widehat b)=\sqrt{\frac{S_1}{\nu r^\top r}},
 \qquad t^2=\frac{\nu D}{S_1}.
\]
Since $S_0=S_1+D$,
\[
 \boxed{R^2_{\rm partial}:=\frac{D}{S_0}
       =\frac{t^2}{t^2+\nu}}.
\]
These algebraic identities require no random model. Under the correctly specified model $y=X\alpha+bz+\varepsilon$, conditional normal errors $\varepsilon\sim N(0,\sigma^2I)$ with $\sigma^2>0$ yield $t\sim t_\nu$ under $b=0$: the normal coefficient component is independent of the orthogonal residual, whose scaled squared norm is $\chi^2_\nu$. Without normality, the finite-sample $t$ law does not follow from these identities.
\insight One extra regressor adds one orthogonal direction. The fitting gain, significance statistic, and partial explained fraction all derive from that projection.

\problem{A03}{Prediction after Sign Selection}{Conditional distributions; prediction risk}
Let $X,Y$ be independent standard normal variables, and retain only $S=\{XY>0\}$. In this conditional population, find the best affine predictor of $Y$ from $X$ and the best predictor under squared loss. Compare their mean squared errors.
\solution For a standard normal variable, sign and magnitude are independent, and $\Prob(S)=1/2$. After selection the signs agree, with each common sign equally likely; the magnitudes remain independent half-normal variables. Thus the conditional means are zero, the variances are one, and
\[
 \E[XY\mid S]=\E|X|\,\E|Y|=\frac2\pi.
\]
The optimal affine predictor is $\widehat Y_{\rm lin}=2X/\pi$, with zero intercept and risk $1-4/\pi^2$.
Given $X=x\ne0$ and $S$, only the sign of $Y$ is determined. Hence
\[
 m(x)=\E[Y\mid X=x,S]
 =\operatorname{sign}(x)\sqrt{2/\pi}.
\]
This minimizes squared loss among all measurable predictors: for any $g$, conditional squared error decomposes into the conditional variance plus $(m(X)-g(X))^2$. The half-normal variance is $1-2/\pi$, so the optimal risk is $1-2/\pi$, improving on the linear risk by
\[
 \frac2\pi-\frac4{\pi^2}>0.
\]
The value assigned at $x=0$ is immaterial because that event has probability zero.
\insight Selection can create correlation and a nonlinear conditional mean. Determine the conditional distribution before comparing linear and unrestricted prediction.

\par\Needspace{8\baselineskip}
\subsection*{Section Exercises}
\exercise{OA02}{Regression on a Residual}
Let $Y,X_1,X_2$ have finite second moments. Assume that both predictors have positive variance and correlation $\rho$ with $|\rho|<1$. Every regression below is a population least-squares regression with an intercept. Regress $Y$ on $X_1$, and denote the residual by $e$. Let $\beta_1$ be the slope from regressing $e$ on $X_2$, and let $\beta_2$ be the coefficient of $X_2$ in the joint regression of $Y$ on $(X_1,X_2)$. Find their relationship.

\exercise{OA03}{Regression after Truncation}
Let $X$ and $Y$ be independent $U[0,1]$ variables. In the conditional population defined by $X+Y>1$, find the slope in the linear regression of $Y$ on $X$ with an intercept.


\section{Estimation}
\noindent\textbf{Key Ideas and Methods.}\enspace
Specify the comparison class before optimizing: linear unbiased estimation, constrained least squares, and posterior estimation use different criteria. With a positive definite covariance matrix, multiplier equations determine optimal linear weights. Completing a quadratic loss expresses the cost of a constraint as a distance in the appropriate metric. Correlated readings cannot be counted as independent information.

An absolute-value penalty requires sign cases or a subgradient at zero. For collinear regressors, separate their combined fitted contribution from its allocation among coefficients, and check uniqueness separately. Distinguish a hard constraint from a finite penalty. For Gaussian observations, separate signal-bearing directions from noise-only directions before conditioning; joint normality justifies the step from zero covariance to independence.
\par\smallskip



\problem{A04}{A Mean from Correlated Observations}{Generalized least squares; optimal weights}
Observe $Y_i=\mu+\varepsilon_i$, $i=1,\ldots,n$, where $n\ge2$, $\E\varepsilon_i=0$, and
\[
 \Cov(\varepsilon_i,\varepsilon_j)=\sigma^2\rho^{|i-j|},
 \qquad \sigma^2>0,\quad0\le\rho<1.
\]
Find the minimum-variance linear unbiased estimator and its variance. Compare it with the arithmetic mean and determine the limit as $\rho\uparrow1$.
\solution Write the correlation matrix as $R$. Direct multiplication verifies $R^{-1}=T/(1-\rho^2)$, where $T$ is tridiagonal: its two endpoint diagonal entries are $1$, its interior diagonal entries are $1+\rho^2$, and adjacent off-diagonal entries are $-\rho$. There are no interior entries when $n=2$.
An estimator $w^\top Y$ is unbiased exactly when $\mathbf1^\top w=1$. Lagrange multipliers minimize $w^\top Rw$ at $w=R^{-1}\mathbf1/(\mathbf1^\top R^{-1}\mathbf1)$. The endpoint row sums are $1/(1+\rho)$; interior row sums are $(1-\rho)/(1+\rho)$. Thus, with $D=2+(n-2)(1-\rho)$,
\[
 \widehat\mu=\frac{Y_1+Y_n+(1-\rho)\sum_{i=2}^{n-1}Y_i}{D},
 \qquad
 \boxed{\Var(\widehat\mu)=\frac{\sigma^2(1+\rho)}{D}}.
\]
Positive definiteness makes the solution unique. Alternatively, any other unbiased weights are $w+h$, where $\mathbf1^\top h=0$. The cross term vanishes, and the variance increases by $\sigma^2h^\top Rh\ge0$.
For $\rho=0$ this is the arithmetic mean. For $n>2$ and $\rho>0$, endpoints receive more weight because interior observations correlate with neighbors on both sides. For $n=2$ the weights remain equal. With $n$ fixed and $\rho\uparrow1$, the endpoint weights approach $1/2$ each and the minimum variance approaches $\sigma^2$, rather than zero.
\insight More readings do not necessarily mean more independent information. Optimal linear weights depend on the inverse covariance matrix; normality is unnecessary.

\problem{A05}{Penalizing Identical Regressors}{Elastic net; uniqueness}
Let $\|x\|^2=1$ and $z=x^\top y$, and suppose both regressors equal $x$. For $\lambda>0$ and $\eta\ge0$, find all minimizers of
\[
 \begin{aligned}
 \min_{b_1,b_2}\quad&\frac12\|y-x(b_1+b_2)\|^2\\
 &+\lambda(|b_1|+|b_2|)+\frac\eta2(b_1^2+b_2^2).
 \end{aligned}
\]
Explain why an arbitrarily small positive quadratic penalty changes uniqueness of the coefficient allocation.
\solution Set $s=b_1+b_2$. For fixed $s$, the triangle inequality gives $|b_1|+|b_2|\ge|s|$, while $b_1^2+b_2^2\ge s^2/2$. Both lower bounds hold with equality at $b_1=b_2=s/2$. Ignoring a constant, the remaining problem is
\[
 \min_s\ \frac12(s-z)^2+\lambda|s|+\frac\eta4s^2.
\]
Differentiating on each side of zero and checking the subgradient at zero yields
\[
 s^*=\frac{\operatorname{sign}(z)(|z|-\lambda)_+}{1+\eta/2}.
\]
If $\eta>0$, equality in the squared-norm bound forces the unique allocation $b_1=b_2=s^*/2$. If $\eta=0$, the sum remains unique. For $s^*>0$, all solutions are $(t,s^*-t)$ with $0\le t\le s^*$; for $s^*<0$, they have $s^*\le t\le0$. If $s^*=0$, the assumption $\lambda>0$ forces both coefficients to be zero.
Thus a nonzero lasso fit admits multiple allocations between identical regressors. Any $\eta>0$ selects the equal allocation and also shrinks their total coefficient. The boundary $|z|=\lambda$ belongs to the zero-solution case.
\insight First fix the combined fitted contribution of collinear variables, then optimize how that contribution is allocated.

\problem{A06}{Least Squares under a Constraint}{Constrained projection; penalty limit}
Let $X$ have full column rank, $A=X^\top X$, and $\widehat\beta=A^{-1}X^\top y$. For a nonzero vector $c$ and a scalar $d$, minimize residual sum of squares subject to $c^\top\beta=d$, and find its increase over the unconstrained minimum. For $\lambda\ge0$, also solve the penalized problem with objective $\|y-X\beta\|^2+\lambda(c^\top\beta-d)^2$ and analyze $\lambda\to\infty$.
\solution The normal equations eliminate the cross term, so
\[
 \|y-X\beta\|^2=\|y-X\widehat\beta\|^2
 +(\beta-\widehat\beta)^\top A(\beta-\widehat\beta).
\]
Write $q=c^\top A^{-1}c>0$ and $e=c^\top\widehat\beta-d$. The multiplier equation puts the correction in direction $A^{-1}c$. Enforcing the constraint gives
\[
 \boxed{\widehat\beta_c=\widehat\beta-\frac e q A^{-1}c},
 \qquad \Delta\mathrm{RSS}=\frac{e^2}{q}.
\]
Positive definiteness guarantees uniqueness. If $e=0$, the constraint costs nothing.
For the penalized minimizer $\beta_\lambda$, the first-order condition is
\[
 A(\beta_\lambda-\widehat\beta)
 +\lambda c(c^\top\beta_\lambda-d)=0.
\]
Multiplication by $c^\top A^{-1}$ shows that the constraint error is $e/(1+\lambda q)$. Therefore
\[
 \beta_\lambda=\widehat\beta-
 \frac{\lambda e}{1+\lambda q}A^{-1}c
 \longrightarrow\widehat\beta_c.
\]
A finite penalty generally does not enforce the constraint exactly. Unless $e=0$, the discrepancy vanishes only at order $1/\lambda$.
\insight This projection uses the metric defined by $A$, not the ordinary Euclidean metric. The initial constraint error and the cost of its correction are measured by $e$ and $q$.

\problem{P34}{Combining Correlated Sensors}{Gaussian conditioning; sufficient statistics}
Let $X\sim N(\mu,s^2)$ with $s^2>0$. Two sensors report $Z_1=X+\varepsilon_1$ and $Z_2=X+\varepsilon_2$. The noise pair is independent of $X$, jointly normal and centered, with common variance $\tau^2>0$ and correlation $-1<\rho<1$. Find the posterior law of $X$, show that the difference of the readings adds no information, and examine both endpoint limits of $\rho$.
\solution Set $M=(Z_1+Z_2)/2$ and $D=Z_1-Z_2$. Then $M=X+\bar\varepsilon$, with
\[
 v:=\Var(\bar\varepsilon)=\frac{\tau^2(1+\rho)}2,
 \quad \Var(D)=2\tau^2(1-\rho).
\]
Equal noise variances give $\Cov(\bar\varepsilon,D)=0$, and $\Cov(X,D)=0$. Joint normality therefore makes $D$ independent of $(X,M)$. Since $(M,D)$ and $(Z_1,Z_2)$ determine each other, the posterior depends only on the average $M$.
Multiplying the prior density by the likelihood of $M$ gives the following part of the exponent involving $x$:
\[
 -\frac12\left\{\frac{(x-\mu)^2}{s^2}
                  +\frac{(M-x)^2}{v}\right\}.
\]
Completing the square yields
\[
 X\mid Z_1,Z_2\sim N(m_*,V_*),\qquad
 V_*=(s^{-2}+v^{-1})^{-1},\quad
 m_*=V_*\left(\frac\mu{s^2}+\frac Mv\right).
\]
As $\rho\uparrow1$, $v\to\tau^2$: the readings become duplicates and carry the information of one sensor. As $\rho\downarrow-1$, $v\to0$, so the average noise cancels and posterior variance tends to zero. In the endpoint model, $M=X$ almost surely. For fixed $s^2,\tau^2$, $V_*$ is strictly increasing in $\rho$.
\insight Decompose multiple measurements into a signal-bearing direction and a pure-noise direction before conditioning. Correlation determines whether the measurements complement or duplicate one another.


\par\Needspace{8\baselineskip}
\subsection*{Section Exercises}
\exercise{OA05}{When Lasso Selects Zero}
For $X\in\R^{n\times p}$ and $y\in\R^n$, with $n\ge1$, consider
\[
 L(\beta)=\frac1{2n}\|y-X\beta\|^2+\lambda\|\beta\|_1,
 \qquad \lambda\ge0.
\]
There is no separate unpenalized intercept. Find a necessary and sufficient condition for $\beta=0$ to be a minimizer.


\section{Matrices and Quadratic Forms}
\noindent\textbf{Key Ideas and Methods.}\enspace
Correlation and Gram matrices are positive semidefinite. Trace, rank, and Rayleigh quotients connect correlations, vector geometry, and eigenvalues. Valid individual entries need not form a compatible matrix. After proving a bound, construct an equality case and allow singular boundary matrices.

With only moment information, optimize a shifted squared-moment bound and check its equality distribution against the event being studied. A strict event may have only a supremum. Conditioning on a sufficient statistic does not increase variance; completeness can then establish optimality within the unbiased class.

For testing simple hypotheses, the likelihood ratio orders outcomes, while the null distribution calibrates size. Normal shrinkage compares total squared-error risk. Check its dimension and covariance assumptions carefully: improvement in total risk does not imply improvement in every coordinate.
\par\smallskip



\problem{P07}{A Missing Correlation}{Correlation completion; construction}
Let $X,Y,Z$ be standardized random variables with $\Corr(X,Y)=a$ and $\Corr(X,Z)=b$, where $|a|,|b|<1$. Find all feasible values of $c=\Corr(Y,Z)$ and construct a joint distribution for each. Specialize to $a=3/5,b=4/5$ and explain the endpoint degeneracies.
\solution Put $U=Y-aX$ and $V=Z-bX$. These residuals are uncorrelated with $X$, and
\[
 \Var(U)=1-a^2,\quad\Var(V)=1-b^2,
 \quad\Cov(U,V)=c-ab.
\]
Cauchy--Schwarz gives the necessary condition
\[
 \boxed{|c-ab|\le\sqrt{(1-a^2)(1-b^2)}}.
\]
Conversely, set $q=(c-ab)/\sqrt{(1-a^2)(1-b^2)}\in[-1,1]$. For independent standard normals $G_0,G_1,G_2$, define
\[
 \begin{aligned}
 X&=G_0,\\
 Y&=aG_0+\sqrt{1-a^2}G_1,\\
 Z&=bG_0+\sqrt{1-b^2}\bigl(qG_1+\sqrt{1-q^2}G_2\bigr).
 \end{aligned}
\]
Direct calculation gives unit variances and the required correlations, proving sufficiency. For $|q|<1$ the transformation is invertible and the correlation matrix is positive definite. For $q=\pm1$, the two residuals are linearly dependent and the matrix is singular.
In the numerical case, both $ab$ and the square-root term equal $12/25$, so $0\le c\le24/25$. At the lower endpoint $V=-3U/4$; at the upper endpoint $V=3U/4$. These are valid correlations: an endpoint introduces an exact linear relation among the variables.
\insight Removing the common variable turns the missing correlation into a residual covariance. The bound and the attaining construction then use the same coordinates.


\problem{A25}{A Sharp One-Sided Tail Bound}{Moment constraints; equality}
Suppose $\E X=0$ and $\Var(X)=\sigma^2>0$. For $t>0$, maximize $\Prob(X\ge t)$ over all such distributions and characterize equality. Is the same bound attained for the strict event $X>t$?
\solution For every $a\ge0$, $X\ge t$ implies $(X+a)^2\ge(t+a)^2$. Therefore
\[
 \Prob(X\ge t)\le\frac{\E(X+a)^2}{(t+a)^2}
 =\frac{\sigma^2+a^2}{(t+a)^2}.
\]
The derivative with respect to $a$ has the sign of $at-\sigma^2$, so the sharpest bound uses $a=\sigma^2/t$:
\[
 \boxed{\Prob(X\ge t)\le\frac{\sigma^2}{\sigma^2+t^2}}.
\]
For equality, the squared term must vanish off the event and equal the threshold on it. Thus $X$ must take only the values $t$ and $-\sigma^2/t$. The mean condition forces $\Prob(X=t)=\sigma^2/(\sigma^2+t^2)$, with the remaining probability at the lower point. Its second moment is $\sigma^2$, so this is the unique equality distribution.
The strict event has the same upper bound. If it attained the bound, the non-strict bound would also be an equality, forcing the distribution just found. But its mass at $t$ does not belong to the strict event, a contradiction. There is consequently no maximum. Using upper point $t+\varepsilon$ and lower point $-\sigma^2/(t+\varepsilon)$ gives strict-tail probabilities approaching the bound, so it is the supremum.
\insight Shift before applying a squared-moment bound to retain one-sided information; optimize the shift as well.

\problem{A29}{Coherence of Unit Vectors}{Gram matrices; trace inequalities}
Let $m>d\ge1$, and let $u_1,\ldots,u_m$ be unit vectors in $\R^d$. Prove
\[
 \max_{i\ne j}|u_i^\top u_j|^2\ge\frac{m-d}{d(m-1)}.
\]
Give two necessary and sufficient conditions for equality, and construct an equality example for $d=2,m=3$.
\solution The Gram matrix $G$ is positive semidefinite, has rank at most $d$, and satisfies $\operatorname{tr}G=m$. Pad its eigenvalues with zeros to obtain $d$ entries. Cauchy--Schwarz yields
\[
 \operatorname{tr}(G^2)\ge\frac{m^2}{d}.
\]
If $M=\max_{i\ne j}|u_i^\top u_j|^2$, then
\[
 \operatorname{tr}(G^2)=m+\sum_{i\ne j}|u_i^\top u_j|^2
 \le m+m(m-1)M.
\]
Combining the inequalities proves the bound. Equality holds exactly when $G$ has $d$ nonzero eigenvalues, all equal to $m/d$, and all off-diagonal inner products have the same absolute value. The first condition is equivalent to
\[
 \sum_{i=1}^m u_i u_i^\top=\frac md I_d.
\]
In the plane, take unit vectors with angles $0,\pi/3,2\pi/3$. Their pairwise inner products have absolute value $1/2$, whose square is $1/4=(3-2)/(2\cdot2)$, attaining the bound.
\insight The spectrum controls the sum of squared inner products, and averaging forces at least one large pair. Absolute size, rather than sign, is the relevant quantity.

\problem{A30}{Eigenvalues after a Rank-One Update}{Minimax principle; spectral perturbation}
A real symmetric matrix $A$ has eigenvalues $\lambda_1\le\cdots\le\lambda_n$. Let $B=A+tuu^\top$, where $t\ge0$ and $u\in\R^n$, and write its eigenvalues as $\mu_1\le\cdots\le\mu_n$. Prove, for $i<n$,
\[
 \lambda_i\le\mu_i\le\lambda_{i+1}.
\]
Also show that $\mu_n\ge\lambda_n$, that each eigenvalue increase is at most $t\|u\|^2$, and that their sum is exactly $t\|u\|^2$.
\solution Every Rayleigh quotient increases by $t(u^\top x)^2/\|x\|^2\ge0$, so the minimax principle gives $\mu_i\ge\lambda_i$. The increase in any quotient is at most $t\|u\|^2$; applying the same principle gives $\mu_i\le\lambda_i+t\|u\|^2$.
For the sharper interlacing upper bound, let $E$ span the first $i+1$ eigenvectors of $A$. If $u\ne0$, the space $E\cap u^\perp$ has dimension at least $i$. The update vanishes there, and the Rayleigh quotient of $A$ is at most $\lambda_{i+1}$. Insert an $i$-dimensional subspace of this intersection into the minimax formula for $B$ to obtain $\mu_i\le\lambda_{i+1}$. If $u=0$ or $t=0$, the assertions are immediate.
Finally, taking traces gives
\[
 \sum_{i=1}^n(\mu_i-\lambda_i)
 =\operatorname{tr}(tuu^\top)=t\|u\|^2.
\]
Thus the total spectral displacement is fixed, and no ordered eigenvalue decreases. When $n=1$, the interlacing assertion is empty but the other statements still hold.
\insight A rank-one perturbation vanishes on a subspace of codimension one. This yields finer location control than a norm bound alone.

\problem{CT01}{Unbiased Estimation of a Squared Probability}{Sufficiency; Rao--Blackwell}
Let $X_1,\ldots,X_n$ be independent $\operatorname{Bernoulli}(p)$ variables, where $n\ge2$ and $0<p<1$. Starting with the unbiased estimator $X_1X_2$, construct an unbiased estimator of $p^2$ with no larger variance. Find its exact variance and prove minimum variance among all unbiased estimators.
\solution Put $S=\sum_iX_i$. Given $S=s$, the success locations are uniform over subsets of size $s$, so
\[
 U:=\E[X_1X_2\mid S]=\frac{S(S-1)}{n(n-1)}.
\]
This conditional law does not involve $p$, making $S$ sufficient. Total variance gives $\Var(X_1X_2)=\Var(U)+\E\Var(X_1X_2\mid S)$, so conditioning preserves unbiasedness and cannot increase variance.
Also, $U$ is the average of all $\binom n2$ products $X_iX_j$. Each has variance $p^2(1-p^2)$; two products sharing one index have covariance $p^3(1-p)$; disjoint products are independent. There are $n\binom{n-1}{2}$ unordered pairs sharing an index. Hence
\[
 \Var(U)=\frac{2p^2(1-p^2)+4(n-2)p^3(1-p)}{n(n-1)}.
\]
To prove minimum variance, first establish completeness. If $\E_p h(S)=0$ for every $p\in(0,1)$, divide by $(1-p)^n$. The resulting polynomial in $p/(1-p)$ vanishes identically, so every coefficient $\binom ns h(s)$ is zero.
For any other unbiased estimator $W$, its conditional mean given $S$ is a parameter-free function with the same expectation as $U$. Completeness forces equality. Total variance then gives $\Var(W)\ge\Var(U)$; an estimator without finite variance cannot improve the bound either. For $n=2$, $U=X_1X_2$ and the variance formula still applies.
\insight Conditioning removes parameter-free randomness. Completeness rules out any different unbiased function of the sufficient statistic that could improve the result.

\problem{CT02}{A Most Powerful Exponential Test}{Likelihood ratios; exact power}
Let $X_1,\ldots,X_n$ be independent exponential variables with rate $\theta$. Test $H_0:\theta=\theta_0$ against $H_1:\theta=\theta_1>\theta_0$. Find the most powerful test among those with type I error at most $\alpha\in(0,1)$, and calculate its power. Give the explicit cutoff and power for $n=2,\theta_0=1,\theta_1=2,\alpha=1-2/e$.
\solution Set $T=\sum_iX_i$. The joint likelihood ratio of alternative to null is
\[
 L=\left(\frac{\theta_1}{\theta_0}\right)^n
      e^{-(\theta_1-\theta_0)T}.
\]
It decreases strictly in $T$, suggesting $\phi_*=\ind_{\{T<c\}}$, where $c>0$ uniquely satisfies $\Prob_{\theta_0}(T<c)=\alpha$. Continuity makes boundary randomization unnecessary.
Let $k$ be the corresponding likelihood-ratio threshold. Every competing test $0\le\phi\le1$ satisfies $(\phi_* -\phi)(L-k)\ge0$. If its size is at most $\alpha$, then
\[
 \E_{\theta_1}(\phi_* -\phi)
 \ge k\E_{\theta_0}(\phi_* -\phi)\ge0.
\]
Thus the candidate maximizes power, proving the likelihood-ratio principle in this setting.
The sum has density $\theta^nt^{n-1}e^{-\theta t}/(n-1)!$, obtainable by induction through convolution. Repeated integration by parts gives
\[
 \Prob_\theta(T<c)=1-e^{-\theta c}
       \sum_{j=0}^{n-1}\frac{(\theta c)^j}{j!}.
\]
Use $\theta_0$ to solve for the cutoff and $\theta_1$ to obtain power. In the specified case, $c=1$ and power is $1-3e^{-2}$. A larger rate makes total waiting time shorter, so rejection lies in the left tail.
\insight The likelihood ratio determines which tail to use; the null distribution calibrates the size. The rejection direction should not be guessed from the parameter's name.

\problem{CT03}{Risk Reduction by Normal Shrinkage}{Stein's identity; squared loss}
Observe $X\sim N_d(\theta,I_d)$, where $d\ge3$ and $\theta\in\R^d$ is unknown. Within the family
\[
 \delta_a(X)=\left(1-\frac a{\|X\|^2}\right)X,
\]
find the constant $a$ minimizing squared-error risk and show that it improves on $X$ for every $\theta$. Define the estimator arbitrarily at $X=0$, a null event. How does the risk change if the same estimator predicts an independent $Y\sim N_d(\theta,I_d)$?
\solution Write $\delta_a=X+g(X)$, where $g(x)=-ax/\|x\|^2$. Integration by parts against the normal density gives
\[
 \E[(X_i-\theta_i)g_i(X)]=\E[\partial_i g_i(X)].
\]
The singularity at zero is harmless here. First integrate over $\|x\|>\varepsilon$: the spherical boundary term has order $\varepsilon^{d-2}\to0$. The derivative has local order $\|x\|^{-2}$, integrable when $d>2$; Gaussian decay removes the boundary at infinity.
Direct differentiation gives
\[
 \|g(x)\|^2=\frac{a^2}{\|x\|^2},\qquad
 \sum_i\partial_i g_i(x)=-\frac{a(d-2)}{\|x\|^2}.
\]
Expanding risk and using the identity yields
\[
 \begin{aligned}
 R_\theta(a)&=\E_\theta\|\delta_a(X)-\theta\|^2\\
 &=d+\bigl(a^2-2a(d-2)\bigr)
                 \E_\theta\frac1{\|X\|^2}.
 \end{aligned}
\]
The final expectation is finite and strictly positive, so the same constant $a=d-2$ minimizes risk for every $\theta$. Its risk is
\[
 d-(d-2)^2\E_\theta\frac1{\|X\|^2}<d,
\]
whereas $X$ has risk $d$. This compares total squared error across coordinates, not the risk of each coordinate separately.
Independence and $\E(Y-\theta)=0$ eliminate the predictive cross term, giving $\E\|Y-\delta_a(X)\|^2=d+R_\theta(a)$. The optimal shrinkage constant is unchanged, and its prediction risk is strictly below the direct predictor's risk $2d$.
\insight A suitable bias can reduce total squared risk. Gaussian integration by parts converts the difficult error cross term into the divergence of the shrinkage adjustment.


\par\Needspace{8\baselineskip}
\subsection*{Section Exercises}
\exercise{OA30}{Eigenvalues of a Principal Submatrix}
Let $A$ be a real symmetric $n\times n$ matrix, with $n\ge2$, and let $B$ be obtained by deleting row $k$ and column $k$. Write the eigenvalues in increasing order as
\[
 \lambda_1\le\cdots\le\lambda_n,
 \qquad \mu_1\le\cdots\le\mu_{n-1}.
\]
Prove the interlacing relation between these eigenvalues.


\section{Exercises}
\noindent\textbf{Key Ideas and Methods.}\enspace
Translate each model into equations and inequalities that can be checked together. For observed moments, test whether one latent model satisfies all covariance restrictions. For probability bounds, separate the upper-bound proof from an attaining construction; mass at the boundary matters for strict events.

For constrained estimation, state feasibility, multiplier signs, complementary slackness, and uniqueness. A shared budget couples the coordinates, so separate thresholding is insufficient. Check the candidate independently by substitution, elimination, or its increase in the objective, and state the parameter range in which the argument applies.
\par\smallskip



\pagebreak
\exercise{AG1-1}{Compatibility of Observed Moments}
Use the independent centered model in A01. The observed moments are
\[
 \begin{aligned}
 \Var(W)&=5,\quad\Var(Z)=4,\quad\Cov(W,Z)=3,\\
 \Cov(W,Y)&=\Cov(Z,Y)=6.
 \end{aligned}
\]
 Identify $\beta$ and the variances of $T,U,V$, and find the regression slope of $Y$ on $W$. If only $\Cov(Z,Y)$ changes to $5$, can the same model still hold?

\exercise{AG1-2}{A Strict Tail Event}
Over all distributions with $\E X=0$ and $\Var(X)=1$, maximize $\Prob(X\ge2)$ and find the supremum of $\Prob(X>2)$. Give the first equality distribution and a sequence approaching the second supremum.

\exercise{AG1-3}{Sparse Estimation with a Budget}
A centered design satisfies $X^\top X/n=I_3$ and $X^\top y/n=(3,1,-1)^\top$. Solve
\[
 \min_{b_j\ge0,\ \sum_jb_j\le2}
 \frac{\|y-Xb\|^2}{2n}+\frac12\sum_{j=1}^3b_j.
\]
Find the budget multiplier and explain why independent coordinate thresholding alone is insufficient.

\exercise{AG1-4}{The Cost of a Coefficient Constraint}
A full-rank least-squares problem has $X^\top X=\operatorname{diag}(1,4)$ and unconstrained solution $(2,-1)^\top$. Impose that the two coefficients sum to zero. Find the constrained solution and the increase in residual sum of squares, and check by direct one-variable minimization.
